\documentclass[10pt,a4paper]{amsart}
\usepackage{amsmath,amssymb,amsthm,graphicx}
\RequirePackage{amssymb,amsmath,amsthm,bm}
\RequirePackage{thmtools}
\RequirePackage[capitalize]{cleveref}
\providecommand{\tightlist}{%
  \setlength{\itemsep}{0pt}\setlength{\parskip}{0pt}}

%%%%% varphi instead of phi
\renewcommand{\phi}{\varphi}

\mathchardef\mhyph="2D 
\providecommand{\A}{\mathcal{A}}
\providecommand{\C}{\mathcal{C}}
\providecommand{\L}{\mathcal{L}}
\providecommand{\B}{\mathcal{B}}
\providecommand{\M}{\mathcal{M}}
\renewcommand{\S}{\mathcal{S}}
\providecommand{\ol}[1]{\overline{#1}}
\providecommand{\mbf}[1]{\mathbf{#1}}
\providecommand{\mf}[1]{\mathfrak{#1}}
\providecommand{\mrm}[1]{\mathrm{#1}}
\providecommand{\mc}[1]{\mathcal{#1}}
\providecommand{\ra}{\rightarrow}
\providecommand{\Ra}{\Rightarrow}
\providecommand{\La}{\Leftarrow}
\providecommand{\lr}{\leftrightarrow}
\providecommand{\hf}{\mathbb{HF}}
\providecommand{\spec}{DgSp}
\providecommand{\bSigma}{\bm\Sigma}
\providecommand{\bDelta}{\bm\Delta}
\providecommand{\bPi}{\bm\Pi}
\providecommand{\Pinf}[1]{\Pi^{\mathrm{in}}_{#1}}
\providecommand{\Sinf}[1]{\Sigma^{\mathrm{in}}_{#1}}
\providecommand{\dSinf}[1]{d\mhyph\Sigma^{\mathrm{in}}_{#1}}
\providecommand{\Dinf}[1]{\Delta^{\mathrm{in}}_{#1}}
\providecommand{\bp}{\bar p}
\providecommand{\bq}{\bar q}
\providecommand{\ba}{\bar a}
\providecommand{\bb}{\bar b}
\providecommand{\bc}{\bar c}
\providecommand{\bd}{\bar d}
\providecommand{\bx}{\bar x}
\providecommand{\by}{\bar y}
\providecommand{\bz}{\bar z}
\providecommand{\LO}{\mathrm{LO}}
\declaretheorem[name={Theorem},numberwithin=section]{theorem}
\declaretheorem[name={Proposition}, sibling=theorem]{proposition}
\declaretheorem[name={Lemma}, sibling=theorem]{lemma}
\declaretheorem[name={Fact}, sibling=theorem]{fact}
\declaretheorem[name={Claim}, numberwithin=theorem]{claim}
\declaretheorem[name={Definition},numberwithin=section,sibling=theorem, style=definition]{definition}
\declaretheorem[name={Remark}, numbered=no,style=definition]{remark}
\declaretheorem[name={Corollary}, sibling=theorem]{corollary}
\declaretheorem[name={Question},sibling=theorem]{question}
\newcommand{\T}{{\mathcal{T}}}
\newcommand{\SR}{SR}
\newcommand{\pSR}{pSR}
\newcommand{\Q}{\mathbb{Q}}
\newcommand{\Z}{\mathbb{Z}}
\newcommand{\K}{\mathcal{K}}
\newcommand{\N}{\mathcal{N}}
\renewcommand{\L}{\mathcal{L}}
\providecommand{\Ginf}[1]{\Gamma^{\mathrm{in}}_{#1}}
\title{Three ordinal-indexed comparisons of products with iterated Z hold at every countable ordinal}
\author{David Gonzalez, Dino Rossegger}
\date{arxiv:2305.07126v1; CC BY 4.0}
\begin{document}
\maketitle
\noindent\textit{Editorial note.} The counted statement is the author environment at source0.tex lines 1147-1158 and its proof environment opens at line 1160, the line immediately after; the pairing is by adjacency and the retained proof range 1160-1266 is the whole environment.\par\medskip
\noindent\textit{Editorial note.} The retained preamble is the author style file dino.sty (package lines, the theorem declaration support, the notation and tuple macros and the theorem environments) followed by the ordinal and rank macros of the manuscript itself, reproduced verbatim; this is the same recipe as the previously published card of this article, so no package outside it is requested and extra\_packages stays empty.\par\medskip
\noindent\textit{Editorial note.} Metadata provenance: source\_id and source\_version are the values recorded in this source licence evidence file (sources/scott-linear-orders-2305.07126v1/LICENSE\_EVIDENCE.json: arxiv:2305.07126v1 / 2305.07126v1), which the card points at directly; the title and author list are transcribed from the authors\_index and title keys of the source record of the pinned copy. The card carries no counter setting and imposes no original numbering, since this article has no label-number dump and its numbering cannot be reconstructed from the retained ranges alone.\par\medskip
\noindent\textit{Editorial note.} Three item labels are defined inside the retained statement, one for each family A\_alpha, B\_alpha and C\_alpha, and the retained proof refers to the families by those names; they need no external entry. No context span and no external reference are needed: the retained statement and proof use no label or citation from outside their own ranges.\par\medskip
\noindent\textit{Editorial scope.} One original result of the pinned author manuscript, counted once: the three ordinal-indexed comparison families for products with iterated Z at the ranks two alpha, two alpha plus one and two alpha plus two (author label lem:timesZ).\par\medskip
\noindent\textit{Allowed standard prerequisites (editorial).}\begin{itemize}
\item Iterated products with Z. Z\textasciicircum{}alpha is the alpha-fold lexicographic product of the integers with itself, and Z\textasciicircum{}alpha times L is the product of that order with L.. They are the orders the three families of statements compare.
\item Back-and-forth relations and their rank shift. The rank two alpha, two alpha plus one and two alpha plus two relations are decided by intervals of the products, whose form fixes how many moves the comparison costs.. They are the relations the three statements assert at those ranks.
\item Initial and final segments of the products. An initial segment of Z\textasciicircum{}alpha times L above a chosen point is again a product with an infinite order, and the same holds for final segments.. They are the orders appearing in the interval cases the proof checks.
\end{itemize}\medskip
\begin{lemma}\label{lem:timesZ}
For the sake of organization, consider the following ordinal indexed propositions.
\begin{enumerate}\tightlist
\item[($A_\alpha$)] For any $\K$ and $\L$, $\Z^\alpha\cdot \K
    \leq_{2\alpha}\Z^\alpha\cdot \L$.
\item[($B_\alpha$)] For any $\K$ and $\L$ with $|\K|\geq|L|$, $\Z^\alpha\cdot \K
    \leq_{2\alpha+1}\Z^\alpha\cdot \L$.
\item[($C_\alpha$)] For any $\K$ and any $\L$ without a last element,
    $\Z^\alpha\cdot(\omega+\K)\leq_{2\alpha+2} \Z^\alpha\cdot(\omega+\L)$.
\end{enumerate}
For any countable ordinal $\alpha$, $A_\alpha$, $B_\alpha$ and $C_\alpha$ are true.
\end{lemma}

\begin{proof}
    The proof is by induction with the statements $A_0$ and $B_0$ trivially true.
    We start by showing that $B_\alpha$ implies $A_{\alpha+1}$. 

    Select a tuple $a_0\cdots a_n$ in $\Z^{\alpha+1}\cdot \L=\Z^\alpha\cdot\Z\cdot
    \L$. We will find a tuple $b_0\dots b_n$ in $\Z^{\alpha+1}\cdot
    \K=\Z^\alpha\cdot\Z\cdot \K$ such that $(\Z^{\alpha+1}\cdot \L,\ba
    )\leq_{2\alpha+1} (\Z^{\alpha+1}\cdot \K,\bb)$. Pick a point $b_0$ arbitrarily.
    Given $b_j$, we pick $b_{j+1}$ as follows. If $a_j$ and $a_{j+1}$ are in the
    same copy of $\Z^{\alpha+1}$, we can pick $b_{j+1}$ such that
    $(a_j,a_{j+1})\cong(b_j,b_{j+1})$. If not, we pick $b_{j+1}$
    arbitrarily in the copy of $\Z^{\alpha}$ that is the successor of the successor
    of the copy of $b_j$. In other words, we pick $b_{j+1}$ with  $(b_j,b_{j+1})\cong \zeta^*_\alpha +\Z^\alpha+\zeta_\alpha$.
    
To show that our choice of $\bb$ is correct, we have to confirm the following
relations between non-isomorphic intervals.
\begin{enumerate}\tightlist
\item $(b_j,b_{j+1})\cong \zeta^*_\alpha +\Z^\alpha+\zeta_\alpha \geq_{2\alpha+1} \zeta^*_\alpha
    +\Z^\alpha\cdot \L_1+\zeta_\alpha$, where $\L_1$ is an infinite order,
\item $(-\infty,b_0)\cong \Z^\alpha\cdot \K_1+\zeta_\alpha \geq_{2\alpha+1}\Z^\alpha\cdot
    \L_1+\zeta_\alpha$, where $\L_1$ and $\K_1$ are some initial segments of the $\Z^\alpha$, and therefore infinite, or
\item $(b_{n},\infty)\cong\zeta^*_\alpha +\Z^\alpha\cdot \L_1 \geq_{2\alpha+1} \zeta^*_\alpha
    +\Z^\alpha\cdot \K_1$, where $\L_1$ and $\K_1$ are some final segments of the $\Z^\alpha$, and therefore infinite.
\end{enumerate} 
It follows immediately from $B_\alpha$ that all of these back-and-forth relations hold.

We now show that $B_\alpha$ implies $C_\alpha$. So assume that $\K$ and $\L$ have
no last element and that $B_\alpha$ holds. 
Select a tuple $a_0\cdots  a_n$ in $\Z^\alpha\cdot(\omega+\L)$.
Without loss of generality, $a_0$ is in the first copy of $\Z^\alpha$. We will
find a tuple $b_0\dots b_n$ in $\Z^\alpha\cdot(\omega+\K)$ such that
$(\Z^\alpha\cdot (\omega+\L),\ba)\leq_{2\alpha+1}
(\Z^\alpha\cdot(\omega+\K),\bb)$. For the $a_j$ in the initial $\Z^\alpha\cdot\omega$, the $b_j$ are picked isomorphically. 
We now define the rest of the $b_j$ by induction so that all of them are in the initial $\Z^\alpha\cdot\omega$.
If $a_j$ is in the same $\Z^\alpha$ block as $a_{j+1}$ or $a_{j+1}$ is in the
succesor of $a_j$'s block, define $b_{j+1}$ so that
$(b_j,b_{j+1})\cong (a_j,a_{j+1})$.
Otherwise, define $b_{j+1}$ as being some element in the successor of the
successor of the $\Z^\alpha$ block of $b_j$. 
Not including isomorphic intervals, we only need to check the following cases. 
\begin{enumerate}
    \item $(b_j,b_{j+1})\cong\zeta^*_\alpha +\Z^\alpha+\zeta_\alpha \geq_{2\alpha+1} \zeta^*_\alpha
    +\Z^\alpha\cdot \L_1+\zeta_\alpha$ where $\L_1$ is infinite, 
\item $(b_n,\infty)\cong\zeta^*_\alpha +\Z^\alpha\cdot (\omega+\K) \geq_{2\alpha+1}
    \zeta^*_\alpha +\Z^\alpha\cdot \L_1$, where $\L_1$ is some final segment of
    $\omega+\L$, and thus infinite.
\end{enumerate} 
All of these relations follow from $B_\alpha$, as desired.

To complete the proof of the successor step of the induction, we show that
$C_\alpha$ and $A_{\alpha+1}$ imply $B_{\alpha+1}$.

Select a tuple  $a_0\cdots a_n$ in $\Z^{\alpha+1}\cdot \L$. These points belong
to at most $n$ distinct copies of $\Z^{\alpha+1}$. Because $|\K|\geq| \L|$, we
can match each chosen copy of $\Z^{\alpha+1}$ to one in $\Z^{\alpha+1}\cdot \K$.
Furthermore, if there is a copy of $\Z^{\alpha+1}$ between two of the chosen
copies in $\Z^{\alpha+1}\cdot \L$, we can assure that the corresponding copies
in $\Z^{\alpha+1}\cdot \K$ are also not successors. Also, if $a_n$ is not in the
last copy of $\Z^{\alpha+1}$, we can assure that neither is $b_n$ and, similarly
for $a_0$ and $b_0$ in the first copy. In particular, not including the
isomorphic intervals, we only need to check the following cases.
\begin{enumerate}
    \item $(b_j,b_{j+1})\cong \zeta^*_{\alpha+1} +\Z^{\alpha+1}\cdot
        \K_1+\zeta_{\alpha+1}\geq_{2\alpha+2}\zeta^*_{\alpha+1}
        +\Z^{\alpha+1}\cdot \L_1+\zeta_{\alpha+1}$ for some orders $\K_1$ and $\L_1$,
    \item $(b_n,\infty)\cong\zeta^*_{\alpha+1} +\Z^{\alpha+1}\cdot
        \K_1\geq_{2\alpha+2}\zeta^*_{\alpha+1} +\Z^{\alpha+1}\cdot \L_1$ for
        some orders $\K_1$ and $\L_1$,
\item $(-\infty,b_0)\cong\Z^{\alpha+1}\cdot
    \K_1+\zeta_{\alpha+1}\geq_{2\alpha+2}\Z^{\alpha+1}\cdot
    \L_1+\zeta_{\alpha+1}$ for some orders $\K_1$ and $\L_1$,
\item $(-\infty,b_0)\cong\zeta_{\alpha+1}\geq_{2\alpha+2} \Z^{\alpha+1}\cdot
    \L_1+\zeta_{\alpha+1}$ for some order $\L_1$,
\item $(b_n,\infty)\cong\zeta^*_{\alpha+1} +\Z^{\alpha+1}\cdot
    \K_1\geq_{2\alpha+2} \zeta^*_{\alpha+1} $ for some order $\K_1$.
\end{enumerate}

The first three cases are handled immediately by $A_{\alpha+1}$. By symmetry it
is sufficient to show case (5). However, this is the same as showing that 
$$\zeta_{\alpha}^* + \Z^\alpha\cdot\omega \leq_{2\alpha+2}\zeta_{\alpha}^* +
\Z^{\alpha}\cdot (\omega+\Z\cdot \K_1).$$
This follows directly from $C_\alpha$.




At last we consider the remaining limit levels. We already have seen that
$B_\lambda$ implies $C_\lambda$. The statement $A_\lambda$ follows immediately from $A_\gamma$ for $\gamma<\lambda$. 
A bit more nuanced is the issue of $B_\lambda$. It follows from $A_\lambda$ and $C_\gamma$ for $\gamma<\lambda$.
Analogous to the proof in the successor case, we analyze the possibilities:
\begin{enumerate}
\item $\zeta^*_{\lambda} +\Z^{\lambda}\cdot
    \K_1+\zeta_{\lambda}\geq_\lambda \zeta^*_{\lambda} +\Z^{\lambda}\cdot \L_1+\zeta_{\lambda} $ for some orders $\L_1$ and $\K_1$,
\item $\zeta^*_{\lambda} +\Z^{\lambda}\cdot \K_1\geq_\lambda \zeta^*_{\lambda} +\Z^{\lambda}\cdot \L_1 $ for some orders $\L_1$ and $\K_1$,
\item $\Z^{\lambda}\cdot \K_1+\zeta_{\lambda}\geq_\lambda \Z^{\lambda}\cdot \L_1+\zeta_{\lambda} $ for some orders $\L_1$ and $\K_1$,
\item $\Z^{\lambda}\cdot \K_1+\zeta_{\lambda}\geq_\lambda \zeta_{\lambda} $ for some order $\K_1$,
\item $\zeta^*_{\lambda} +\Z^{\lambda}\cdot \K_1\geq_\lambda\zeta^*_{\lambda} $ for some order $\K_1$.
\end{enumerate}

The first three cases are handled immediately by $A_{\lambda}$. By symmetry it is clear that it is enough to show case 5. However, this is the same as showing that for all $\gamma<\lambda$,
$$\zeta_{\gamma}^* + \Z^\gamma\cdot(\omega +  \Z\cdot\zeta_{\lambda}^*)
\leq_{\gamma}\zeta_{\gamma}^* + \Z^{\gamma}\cdot
(\Z\cdot\zeta_{\lambda}^*+\Z^\lambda\cdot \K_1),$$
which follows from the $C_\gamma$ for $\gamma<\lambda$.

Therefore, for any countable ordinal $\alpha$, $A_\alpha$, $B_\alpha$ and $C_\alpha$ all hold.
\end{proof}

\end{document}
